% =====================================================================
%  PAGINA DI LAVORO — proposte di titolo
%  Dopo la scelta: aggiornare \title in frontmatter/titlepage.tex e
%  rimuovere la riga % =====================================================================
%  PAGINA DI LAVORO — proposte di titolo
%  Dopo la scelta: aggiornare \title in frontmatter/titlepage.tex e
%  rimuovere la riga % =====================================================================
%  PAGINA DI LAVORO — proposte di titolo
%  Dopo la scelta: aggiornare \title in frontmatter/titlepage.tex e
%  rimuovere la riga % =====================================================================
%  PAGINA DI LAVORO — proposte di titolo
%  Dopo la scelta: aggiornare \title in frontmatter/titlepage.tex e
%  rimuovere la riga \input{frontmatter/proposte-titolo} da main.tex.
% =====================================================================
\thispagestyle{empty}

\begin{fullwidth}
\begin{center}
{\LARGE\scshape Proposte di titolo}\\[2pt]
{\small documento di lavoro --- 21 agosto 2026}
\end{center}

\vspace{0.5em}
\noindent
Elenco di titoli alternativi, raggruppati per taglio. La scelta va concordata con
il relatore.

\vspace{1em}
\noindent\textbf{A. Taglio ``confronto sperimentale'' (il pi\`u aderente al contenuto)}
\begin{enumerate}[label=A\arabic*.,leftmargin=2.2em,itemsep=3pt]
  \item Sensori di presenza per la sicurezza sismica: confronto sperimentale tra
        radar mmWave a 24~GHz e sensori PIR su piattaforma ESP32
  \item Rilevare chi non si muove: valutazione sperimentale di radar mmWave e PIR
        per il rilevamento di presenza in scenari di emergenza
  \item Caratterizzazione e confronto quantitativo di sensori PIR e radar mmWave a
        24~GHz per il rilevamento di persone in ambienti indoor
  \item Oltre il movimento: analisi comparata PIR / mmWave FMCW per il rilevamento
        di presenza statica
\end{enumerate}

\vspace{0.6em}
\noindent\textbf{B. Taglio ``progetto DIPME'' (contesto applicativo in evidenza)}
\begin{enumerate}[label=B\arabic*.,leftmargin=2.2em,itemsep=3pt]
  \item Sensing di presenza per arredi salva-vita: il radar mmWave a 24~GHz come
        evoluzione del nodo PIR nel progetto SAFE
  \item Un sensore per chi \`e rimasto sotto: rilevamento di presenza e di vitalit\`a
        con radar mmWave per il soccorso post-sismico
  \item Dal ``c'\`e qualcuno'' al ``in che condizioni \`e'': sensing mmWave a supporto
        del triage in emergenza sismica
  \item Nodi IoT per arredi antisismici: valutazione di sensori mmWave, PIR e UWB
        per il rilevamento di persone intrappolate
\end{enumerate}

\vspace{0.6em}
\noindent\textbf{C. Taglio ``sistema realizzato'' (firmware, web UI, indice di vitalit\`a)}
\begin{enumerate}[label=C\arabic*.,leftmargin=2.2em,itemsep=3pt]
  \item Progetto e realizzazione di un nodo di sensing mmWave su ESP32 con
        dashboard web e indice di vitalit\`a
  \item Acquisizione, visualizzazione e analisi di dati radar mmWave a 24~GHz: un
        sistema completo su ESP32 per il rilevamento di presenza
  \item Da micro-movimenti a indice di vitalit\`a: elaborazione di dati radar mmWave
        su microcontrollore
\end{enumerate}

\vspace{0.6em}
\noindent\textbf{D. Taglio tecnico e breve}
\begin{enumerate}[label=D\arabic*.,leftmargin=2.2em,itemsep=3pt]
  \item Radar mmWave a 24~GHz per il rilevamento di presenza statica: misure e
        confronto con sensori PIR
  \item Presence sensing con mmWave ed ESP32: caratterizzazione sperimentale
  \item Micro-movimenti e respiro: analisi dei dati per-gate di un radar FMCW a
        24~GHz
\end{enumerate}

\vspace{1.2em}
\noindent\textbf{Nota per la scelta.} Il risultato sperimentale centrale gi\`a
acquisito \`e la \emph{cecit\`a del PIR sulla persona ferma}: falsi negativi al
$99{,}78 \pm 0{,}49\,\%$ contro lo $0{,}00\,\%$ del radar. I titoli A2, B2 e D1 sono
quelli che lo mettono in prima linea. Se il relatore preferisce che il titolo
espliciti il progetto europeo, la coppia consigliata \`e B1 (sobria) oppure B2
(pi\`u comunicativa). Il taglio C \`e appropriato solo se l'implementazione della
web UI e dell'indice di vitalit\`a saranno completate: al momento sono progettate
ma non realizzate.
\end{fullwidth}

\cleardoublepage
 da main.tex.
% =====================================================================
\thispagestyle{empty}

\begin{fullwidth}
\begin{center}
{\LARGE\scshape Proposte di titolo}\\[2pt]
{\small documento di lavoro --- 21 agosto 2026}
\end{center}

\vspace{0.5em}
\noindent
Elenco di titoli alternativi, raggruppati per taglio. La scelta va concordata con
il relatore.

\vspace{1em}
\noindent\textbf{A. Taglio ``confronto sperimentale'' (il pi\`u aderente al contenuto)}
\begin{enumerate}[label=A\arabic*.,leftmargin=2.2em,itemsep=3pt]
  \item Sensori di presenza per la sicurezza sismica: confronto sperimentale tra
        radar mmWave a 24~GHz e sensori PIR su piattaforma ESP32
  \item Rilevare chi non si muove: valutazione sperimentale di radar mmWave e PIR
        per il rilevamento di presenza in scenari di emergenza
  \item Caratterizzazione e confronto quantitativo di sensori PIR e radar mmWave a
        24~GHz per il rilevamento di persone in ambienti indoor
  \item Oltre il movimento: analisi comparata PIR / mmWave FMCW per il rilevamento
        di presenza statica
\end{enumerate}

\vspace{0.6em}
\noindent\textbf{B. Taglio ``progetto DIPME'' (contesto applicativo in evidenza)}
\begin{enumerate}[label=B\arabic*.,leftmargin=2.2em,itemsep=3pt]
  \item Sensing di presenza per arredi salva-vita: il radar mmWave a 24~GHz come
        evoluzione del nodo PIR nel progetto SAFE
  \item Un sensore per chi \`e rimasto sotto: rilevamento di presenza e di vitalit\`a
        con radar mmWave per il soccorso post-sismico
  \item Dal ``c'\`e qualcuno'' al ``in che condizioni \`e'': sensing mmWave a supporto
        del triage in emergenza sismica
  \item Nodi IoT per arredi antisismici: valutazione di sensori mmWave, PIR e UWB
        per il rilevamento di persone intrappolate
\end{enumerate}

\vspace{0.6em}
\noindent\textbf{C. Taglio ``sistema realizzato'' (firmware, web UI, indice di vitalit\`a)}
\begin{enumerate}[label=C\arabic*.,leftmargin=2.2em,itemsep=3pt]
  \item Progetto e realizzazione di un nodo di sensing mmWave su ESP32 con
        dashboard web e indice di vitalit\`a
  \item Acquisizione, visualizzazione e analisi di dati radar mmWave a 24~GHz: un
        sistema completo su ESP32 per il rilevamento di presenza
  \item Da micro-movimenti a indice di vitalit\`a: elaborazione di dati radar mmWave
        su microcontrollore
\end{enumerate}

\vspace{0.6em}
\noindent\textbf{D. Taglio tecnico e breve}
\begin{enumerate}[label=D\arabic*.,leftmargin=2.2em,itemsep=3pt]
  \item Radar mmWave a 24~GHz per il rilevamento di presenza statica: misure e
        confronto con sensori PIR
  \item Presence sensing con mmWave ed ESP32: caratterizzazione sperimentale
  \item Micro-movimenti e respiro: analisi dei dati per-gate di un radar FMCW a
        24~GHz
\end{enumerate}

\vspace{1.2em}
\noindent\textbf{Nota per la scelta.} Il risultato sperimentale centrale gi\`a
acquisito \`e la \emph{cecit\`a del PIR sulla persona ferma}: falsi negativi al
$99{,}78 \pm 0{,}49\,\%$ contro lo $0{,}00\,\%$ del radar. I titoli A2, B2 e D1 sono
quelli che lo mettono in prima linea. Se il relatore preferisce che il titolo
espliciti il progetto europeo, la coppia consigliata \`e B1 (sobria) oppure B2
(pi\`u comunicativa). Il taglio C \`e appropriato solo se l'implementazione della
web UI e dell'indice di vitalit\`a saranno completate: al momento sono progettate
ma non realizzate.
\end{fullwidth}

\cleardoublepage
 da main.tex.
% =====================================================================
\thispagestyle{empty}

\begin{fullwidth}
\begin{center}
{\LARGE\scshape Proposte di titolo}\\[2pt]
{\small documento di lavoro --- 21 agosto 2026}
\end{center}

\vspace{0.5em}
\noindent
Elenco di titoli alternativi, raggruppati per taglio. La scelta va concordata con
il relatore.

\vspace{1em}
\noindent\textbf{A. Taglio ``confronto sperimentale'' (il pi\`u aderente al contenuto)}
\begin{enumerate}[label=A\arabic*.,leftmargin=2.2em,itemsep=3pt]
  \item Sensori di presenza per la sicurezza sismica: confronto sperimentale tra
        radar mmWave a 24~GHz e sensori PIR su piattaforma ESP32
  \item Rilevare chi non si muove: valutazione sperimentale di radar mmWave e PIR
        per il rilevamento di presenza in scenari di emergenza
  \item Caratterizzazione e confronto quantitativo di sensori PIR e radar mmWave a
        24~GHz per il rilevamento di persone in ambienti indoor
  \item Oltre il movimento: analisi comparata PIR / mmWave FMCW per il rilevamento
        di presenza statica
\end{enumerate}

\vspace{0.6em}
\noindent\textbf{B. Taglio ``progetto DIPME'' (contesto applicativo in evidenza)}
\begin{enumerate}[label=B\arabic*.,leftmargin=2.2em,itemsep=3pt]
  \item Sensing di presenza per arredi salva-vita: il radar mmWave a 24~GHz come
        evoluzione del nodo PIR nel progetto SAFE
  \item Un sensore per chi \`e rimasto sotto: rilevamento di presenza e di vitalit\`a
        con radar mmWave per il soccorso post-sismico
  \item Dal ``c'\`e qualcuno'' al ``in che condizioni \`e'': sensing mmWave a supporto
        del triage in emergenza sismica
  \item Nodi IoT per arredi antisismici: valutazione di sensori mmWave, PIR e UWB
        per il rilevamento di persone intrappolate
\end{enumerate}

\vspace{0.6em}
\noindent\textbf{C. Taglio ``sistema realizzato'' (firmware, web UI, indice di vitalit\`a)}
\begin{enumerate}[label=C\arabic*.,leftmargin=2.2em,itemsep=3pt]
  \item Progetto e realizzazione di un nodo di sensing mmWave su ESP32 con
        dashboard web e indice di vitalit\`a
  \item Acquisizione, visualizzazione e analisi di dati radar mmWave a 24~GHz: un
        sistema completo su ESP32 per il rilevamento di presenza
  \item Da micro-movimenti a indice di vitalit\`a: elaborazione di dati radar mmWave
        su microcontrollore
\end{enumerate}

\vspace{0.6em}
\noindent\textbf{D. Taglio tecnico e breve}
\begin{enumerate}[label=D\arabic*.,leftmargin=2.2em,itemsep=3pt]
  \item Radar mmWave a 24~GHz per il rilevamento di presenza statica: misure e
        confronto con sensori PIR
  \item Presence sensing con mmWave ed ESP32: caratterizzazione sperimentale
  \item Micro-movimenti e respiro: analisi dei dati per-gate di un radar FMCW a
        24~GHz
\end{enumerate}

\vspace{1.2em}
\noindent\textbf{Nota per la scelta.} Il risultato sperimentale centrale gi\`a
acquisito \`e la \emph{cecit\`a del PIR sulla persona ferma}: falsi negativi al
$99{,}78 \pm 0{,}49\,\%$ contro lo $0{,}00\,\%$ del radar. I titoli A2, B2 e D1 sono
quelli che lo mettono in prima linea. Se il relatore preferisce che il titolo
espliciti il progetto europeo, la coppia consigliata \`e B1 (sobria) oppure B2
(pi\`u comunicativa). Il taglio C \`e appropriato solo se l'implementazione della
web UI e dell'indice di vitalit\`a saranno completate: al momento sono progettate
ma non realizzate.
\end{fullwidth}

\cleardoublepage
 da main.tex.
% =====================================================================
\thispagestyle{empty}

\begin{fullwidth}
\begin{center}
{\LARGE\scshape Proposte di titolo}\\[2pt]
{\small documento di lavoro --- 21 agosto 2026}
\end{center}

\vspace{0.5em}
\noindent
Elenco di titoli alternativi, raggruppati per taglio. La scelta va concordata con
il relatore.

\vspace{1em}
\noindent\textbf{A. Taglio ``confronto sperimentale'' (il pi\`u aderente al contenuto)}
\begin{enumerate}[label=A\arabic*.,leftmargin=2.2em,itemsep=3pt]
  \item Sensori di presenza per la sicurezza sismica: confronto sperimentale tra
        radar mmWave a 24~GHz e sensori PIR su piattaforma ESP32
  \item Rilevare chi non si muove: valutazione sperimentale di radar mmWave e PIR
        per il rilevamento di presenza in scenari di emergenza
  \item Caratterizzazione e confronto quantitativo di sensori PIR e radar mmWave a
        24~GHz per il rilevamento di persone in ambienti indoor
  \item Oltre il movimento: analisi comparata PIR / mmWave FMCW per il rilevamento
        di presenza statica
\end{enumerate}

\vspace{0.6em}
\noindent\textbf{B. Taglio ``progetto DIPME'' (contesto applicativo in evidenza)}
\begin{enumerate}[label=B\arabic*.,leftmargin=2.2em,itemsep=3pt]
  \item Sensing di presenza per arredi salva-vita: il radar mmWave a 24~GHz come
        evoluzione del nodo PIR nel progetto SAFE
  \item Un sensore per chi \`e rimasto sotto: rilevamento di presenza e di vitalit\`a
        con radar mmWave per il soccorso post-sismico
  \item Dal ``c'\`e qualcuno'' al ``in che condizioni \`e'': sensing mmWave a supporto
        del triage in emergenza sismica
  \item Nodi IoT per arredi antisismici: valutazione di sensori mmWave, PIR e UWB
        per il rilevamento di persone intrappolate
\end{enumerate}

\vspace{0.6em}
\noindent\textbf{C. Taglio ``sistema realizzato'' (firmware, web UI, indice di vitalit\`a)}
\begin{enumerate}[label=C\arabic*.,leftmargin=2.2em,itemsep=3pt]
  \item Progetto e realizzazione di un nodo di sensing mmWave su ESP32 con
        dashboard web e indice di vitalit\`a
  \item Acquisizione, visualizzazione e analisi di dati radar mmWave a 24~GHz: un
        sistema completo su ESP32 per il rilevamento di presenza
  \item Da micro-movimenti a indice di vitalit\`a: elaborazione di dati radar mmWave
        su microcontrollore
\end{enumerate}

\vspace{0.6em}
\noindent\textbf{D. Taglio tecnico e breve}
\begin{enumerate}[label=D\arabic*.,leftmargin=2.2em,itemsep=3pt]
  \item Radar mmWave a 24~GHz per il rilevamento di presenza statica: misure e
        confronto con sensori PIR
  \item Presence sensing con mmWave ed ESP32: caratterizzazione sperimentale
  \item Micro-movimenti e respiro: analisi dei dati per-gate di un radar FMCW a
        24~GHz
\end{enumerate}

\vspace{1.2em}
\noindent\textbf{Nota per la scelta.} Il risultato sperimentale centrale gi\`a
acquisito \`e la \emph{cecit\`a del PIR sulla persona ferma}: falsi negativi al
$99{,}78 \pm 0{,}49\,\%$ contro lo $0{,}00\,\%$ del radar. I titoli A2, B2 e D1 sono
quelli che lo mettono in prima linea. Se il relatore preferisce che il titolo
espliciti il progetto europeo, la coppia consigliata \`e B1 (sobria) oppure B2
(pi\`u comunicativa). Il taglio C \`e appropriato solo se l'implementazione della
web UI e dell'indice di vitalit\`a saranno completate: al momento sono progettate
ma non realizzate.
\end{fullwidth}

\cleardoublepage
